% !TEX root = ../main.tex
\chapter*{Conclusion générale et perspectives}
\addcontentsline{toc}{chapter}{Conclusion générale et perspectives}
\markboth{Conclusion générale et perspectives}{}

% Plan conseillé : rappel du problème -> ce qui a été réalisé -> apports personnels -> perspectives.

Ce projet de fin d'études, réalisé au sein de \entreprise, avait pour objectif de concevoir et de réaliser \nomprojet, une plateforme SaaS de prise de rendez-vous en ligne pour les cliniques.

Après l'étude de l'existant et des solutions du marché, nous avons spécifié les besoins, conçu une architecture multi-cliniques fondée sur Spring Boot, React et PostgreSQL, puis réalisé la solution en trois releases. La plateforme permet aujourd'hui aux patients de réserver en ligne, aux secrétariats de gérer les agendas et aux cliniques de suivre leur activité. \todo{Résumer les résultats chiffrés.}

Ce stage nous a permis d'approfondir nos compétences en architecture logicielle, en sécurité des applications web et en DevOps, et de travailler au sein d'une équipe agile.

Plusieurs perspectives s'ouvrent à ce travail : la téléconsultation par vidéo, une application mobile pour les patients, et l'intégration avec les dossiers patients informatisés des cliniques.
